\section{Analytic Number Theory}

This section develops the analytic tools used to study arithmetic functions and the distribution of primes. Its main themes are Dirichlet series and their Euler products, Dirichlet characters and $L$-functions, the zeta functions, and the von Mangoldt function. The organization follows Apostol's treatment of arithmetical functions, Dirichlet series, the functions $\zeta(s)$ and $L(s,\chi)$, and the analytic proof of the prime number theorem \cite[Chapters 2--4, 6--7, and 11--13]{Apostol1976-xj}.

\subsection{Arithmetic Functions and Dirichlet Convolution}

\begin{definition}[Arithmetic function]
An arithmetic function is a function defined on the positive integers. For arithmetic functions $f$ and $g$, their Dirichlet convolution is
\[
 (f*g)(n)=\sum_{d\mid n}f(d)g\!\left(\frac{n}{d}\right).
\]
The identity for this operation is the function $\varepsilon$ given by $\varepsilon(1)=1$ and $\varepsilon(n)=0$ for $n>1$.
\end{definition}

\begin{theorem}[Multiplication of Dirichlet series]
Suppose that the Dirichlet series
\[
 F(s)=\sum_{n=1}^{\infty}\frac{f(n)}{n^s},
 \qquad
 G(s)=\sum_{n=1}^{\infty}\frac{g(n)}{n^s}
\]
converge absolutely in a common half-plane. Then their product there is the Dirichlet series
\[
 F(s)G(s)=\sum_{n=1}^{\infty}\frac{(f*g)(n)}{n^s}.
\]
\end{theorem}

\begin{proof}
Absolute convergence permits multiplication and regrouping by the product of the indices:
\[
 F(s)G(s)=\sum_{m=1}^{\infty}\sum_{n=1}^{\infty}
 \frac{f(m)g(n)}{(mn)^s}
 =\sum_{k=1}^{\infty}\frac{1}{k^s}\sum_{d\mid k}f(d)g\!\left(\frac{k}{d}\right).
\]
The inner sum is $(f*g)(k)$, as required.
\end{proof}

\begin{example}[The square of the zeta series]
Let $d(n)$ denote the number of positive divisors of $n$. Since $(1*1)(n)=d(n)$, for $\operatorname{Re}(s)>1$ we have
\[
 \zeta(s)^2=\left(\sum_{n=1}^{\infty}\frac{1}{n^s}\right)^2
 =\sum_{n=1}^{\infty}\frac{d(n)}{n^s}.
\]
For instance, the coefficient of $6^{-s}$ is $d(6)=4$, corresponding to the divisors $1,2,3,6$.
\end{example}

\subsection{The von Mangoldt Function}

\begin{definition}[von Mangoldt function]
For a positive integer $n$, define
\[
 \Lambda(n)=
 \begin{cases}
 \log p,& n=p^k\text{ for a prime }p\text{ and an integer }k\geq1,\\
 0,&\text{otherwise}.
 \end{cases}
\]
In particular, $\Lambda(1)=0$.
\end{definition}

\begin{theorem}[von Mangoldt identity]
For every positive integer $n$,
\[
 \log n=\sum_{d\mid n}\Lambda(d).
\]
\end{theorem}

\begin{proof}
The identity is immediate for $n=1$. If $n=\prod_{j=1}^r p_j^{\alpha_j}$ is the prime factorization of $n>1$, the only divisors $d$ for which $\Lambda(d)\neq0$ are $p_j,p_j^2,\ldots,p_j^{\alpha_j}$. Thus
\[
 \sum_{d\mid n}\Lambda(d)
 =\sum_{j=1}^r\alpha_j\log p_j
 =\log\left(\prod_{j=1}^r p_j^{\alpha_j}\right)=\log n.
\]
\end{proof}

\begin{example}
For $n=360=2^3\cdot3^2\cdot5$, the prime-power divisors contribute
\[
 \sum_{d\mid360}\Lambda(d)=3\log2+2\log3+\log5=\log360.
\]
\end{example}

\subsection{Dirichlet Series and Euler Products}

\begin{definition}[Dirichlet series]
A Dirichlet series with coefficients $a(n)$ is a series of the form
\[
 F(s)=\sum_{n=1}^{\infty}\frac{a(n)}{n^s},
 \qquad s=\sigma+it\in\mathbb{C},
\]
where $n^{-s}=e^{-s\log n}$. Its half-plane of absolute convergence is the set of $s$ for which $\sum_{n\geq1}|a(n)|n^{-\operatorname{Re}(s)}$ converges.
\end{definition}

\begin{theorem}[Abscissa of absolute convergence]
If a Dirichlet series converges absolutely at some real value of $s$, there is an extended real number $\sigma_a$ such that it converges absolutely for $\operatorname{Re}(s)>\sigma_a$ and does not converge absolutely for $\operatorname{Re}(s)<\sigma_a$. No general conclusion is implied on the boundary line $\operatorname{Re}(s)=\sigma_a$.
\end{theorem}

\begin{theorem}[Holomorphy in the convergence half-plane]
A Dirichlet series converges uniformly on compact subsets of its half-plane of absolute convergence and defines a holomorphic function there. In particular, it may be differentiated term by term in that half-plane.
\end{theorem}

\begin{theorem}[Euler product]
Suppose $a$ is multiplicative and $\sum_{n\geq1}a(n)n^{-s}$ converges absolutely. Then
\[
 \sum_{n=1}^{\infty}\frac{a(n)}{n^s}
 =\prod_{p}\left(\sum_{k=0}^{\infty}\frac{a(p^k)}{p^{ks}}\right),
\]
where the product is over all primes. If $a$ is completely multiplicative, this becomes
\[
 \sum_{n=1}^{\infty}\frac{a(n)}{n^s}
 =\prod_p\left(1-\frac{a(p)}{p^s}\right)^{-1}.
\]
\end{theorem}

\begin{proof}
For a finite set of primes, expand the product of the local series. Multiplicativity identifies each resulting coefficient with $a(n)$ for integers whose prime factors lie in that set. Absolute convergence allows the finite set of primes to increase to all primes and the terms to be regrouped, giving the stated identity. In the completely multiplicative case, each local series is geometric.
\end{proof}

\begin{example}[Euler product for $\zeta(s)$]
For $\operatorname{Re}(s)>1$, unique factorization and the geometric series give
\[
 \zeta(s)=\sum_{n=1}^{\infty}\frac{1}{n^s}
 =\prod_p\left(1-p^{-s}\right)^{-1}.
\]
For example, at $s=2$ this expresses the convergent series $\sum_{n\geq1}n^{-2}$ as a product over primes.
\end{example}

Taking logarithms of the Euler product, with absolute convergence for $\operatorname{Re}(s)>1$, yields
\[
 \log\zeta(s)=\sum_p\sum_{k=1}^{\infty}\frac{1}{k p^{ks}}
 =\sum_{n=2}^{\infty}\frac{\Lambda(n)}{\log n}\,\frac{1}{n^s}.
\]
Differentiation gives the fundamental identity
\[
 -\frac{\zeta'(s)}{\zeta(s)}
 =\sum_{p}\sum_{k=1}^{\infty}\frac{\log p}{p^{ks}}
 =\sum_{n=1}^{\infty}\frac{\Lambda(n)}{n^s},
 \qquad \operatorname{Re}(s)>1.
\]
The logarithmic derivative therefore records prime powers through $\Lambda(n)$.

\subsection{Dirichlet Characters and \texorpdfstring{$L$}{L}-Functions}

\begin{definition}[Dirichlet character]
A Dirichlet character modulo $q$ is a periodic, completely multiplicative arithmetic function $\chi$ such that $\chi(n)=0$ when $\gcd(n,q)>1$ and $|\chi(n)|=1$ when $\gcd(n,q)=1$. The character that equals $1$ on integers coprime to $q$ is the principal character modulo $q$.
\end{definition}

\begin{proposition}[Sum of a nonprincipal character]
If $\chi$ is a nonprincipal Dirichlet character modulo $q$, then
\[
 \sum_{a=1}^{q}\chi(a)=0.
\]
\end{proposition}

\begin{proof}
Choose an integer $b$ coprime to $q$ for which $\chi(b)\neq1$. Multiplication by $b$ permutes the residue classes coprime to $q$. Therefore
\[
 \sum_{a=1}^{q}\chi(a)=\sum_{a=1}^{q}\chi(ba)
 =\chi(b)\sum_{a=1}^{q}\chi(a),
\]
which forces the sum to be zero.
\end{proof}

\begin{definition}[Dirichlet $L$-function]
For a Dirichlet character $\chi$ modulo $q$, its Dirichlet $L$-function is
\[
 L(s,\chi)=\sum_{n=1}^{\infty}\frac{\chi(n)}{n^s},
 \qquad \operatorname{Re}(s)>1.
\]
\end{definition}

\begin{theorem}[Euler product for Dirichlet $L$-functions]
For $\operatorname{Re}(s)>1$,
\[
 L(s,\chi)=\prod_p\left(1-\frac{\chi(p)}{p^s}\right)^{-1}.
\]
\end{theorem}

\begin{proof}
A Dirichlet character is completely multiplicative, so the Euler product theorem applies. If $p\mid q$, then $\chi(p)=0$ and the corresponding factor is $1$.
\end{proof}

\begin{example}[The Dirichlet beta function]
The nonprincipal character modulo $4$ is given by $\chi_4(n)=0$ for even $n$, $\chi_4(n)=1$ when $n\equiv1\pmod4$, and $\chi_4(n)=-1$ when $n\equiv3\pmod4$. Its $L$-function is the Dirichlet beta function:
\[
 \beta(s)=L(s,\chi_4)
 =1-\frac{1}{3^s}+\frac{1}{5^s}-\frac{1}{7^s}+\cdots,
 \qquad \operatorname{Re}(s)>1.
\]
At $s=1$, the alternating series gives $\beta(1)=\pi/4$.
\end{example}

\begin{theorem}[Dirichlet's theorem on primes in arithmetic progressions]
If $\gcd(a,q)=1$, then there are infinitely many primes congruent to $a$ modulo $q$.
\end{theorem}

Dirichlet characters separate residue classes, and the corresponding $L$-functions encode the primes in those classes in their Euler products. In Apostol's proof, the theorem relies in particular on the nonvanishing of $L(1,\chi)$ for nonprincipal characters.

\subsection{The Riemann, Eta, and Hurwitz Zeta Functions}

\begin{definition}[Riemann zeta function]
For $\operatorname{Re}(s)>1$, the Riemann zeta function is
\[
 \zeta(s)=\sum_{n=1}^{\infty}\frac{1}{n^s}.
\]
Its Euler product is $\zeta(s)=\prod_p(1-p^{-s})^{-1}$ in the same half-plane.
\end{definition}

\begin{definition}[Dirichlet eta function]
For $\operatorname{Re}(s)>0$, the Dirichlet eta function is
\[
 \eta(s)=\sum_{n=1}^{\infty}\frac{(-1)^{n-1}}{n^s}.
\]
\end{definition}

\begin{proposition}[Relation between $\eta$ and $\zeta$]
For $\operatorname{Re}(s)>1$,
\[
 \eta(s)=\left(1-2^{1-s}\right)\zeta(s).
\]
The eta series is holomorphic for $\operatorname{Re}(s)>0$, so $\zeta(s)=\eta(s)/(1-2^{1-s})$ continues $\zeta$ to that half-plane wherever the denominator is nonzero; values at its zeros are understood by meromorphic continuation.
\end{proposition}

\begin{proof}
Subtract twice the sum over even integers from the full zeta series:
\[
 \eta(s)=\sum_{n=1}^{\infty}n^{-s}-2\sum_{n=1}^{\infty}(2n)^{-s}
 =\zeta(s)-2^{1-s}\zeta(s).
\]
\end{proof}

\begin{example}
At $s=1$, the alternating harmonic series converges and gives
\[
 \eta(1)=1-\frac12+\frac13-\frac14+\cdots=\log2.
\]
The relation above is initially proved for $\operatorname{Re}(s)>1$; using the continuation of $\zeta$, it is not interpreted at $s=1$ by direct substitution into the product $\left(1-2^{1-s}\right)\zeta(s)$ without taking a limit.
\end{example}

\begin{definition}[Hurwitz zeta function]
For $a>0$ and $\operatorname{Re}(s)>1$, the Hurwitz zeta function is
\[
 \zeta(s,a)=\sum_{n=0}^{\infty}\frac{1}{(n+a)^s}.
\]
In particular, $\zeta(s,1)=\zeta(s)$.
\end{definition}

\begin{proposition}[Hurwitz representation of Dirichlet $L$-functions]
For $\operatorname{Re}(s)>1$,
\[
 L(s,\chi)=\frac{1}{q^s}\sum_{a=1}^{q}\chi(a)\zeta\!\left(s,\frac{a}{q}\right).
\]
\end{proposition}

\begin{proof}
Every positive integer has a unique expression $qn+a$ with $n\geq0$ and $1\leq a\leq q$. Grouping the absolutely convergent series for $L(s,\chi)$ by these residue classes gives
\[
 L(s,\chi)=\sum_{a=1}^{q}\chi(a)\sum_{n=0}^{\infty}\frac{1}{(qn+a)^s}
 =\frac{1}{q^s}\sum_{a=1}^{q}\chi(a)\zeta\!\left(s,\frac{a}{q}\right).
\]
\end{proof}

\begin{example}[A special value of the Hurwitz zeta function]
For $\operatorname{Re}(s)>1$, separating the odd terms in the zeta series gives
\[
 \zeta\!\left(s,\frac12\right)
 =2^s\sum_{n=0}^{\infty}\frac{1}{(2n+1)^s}
 =(2^s-1)\zeta(s).
\]
\end{example}

\begin{theorem}[Meromorphic continuation of the Hurwitz zeta function]
For each $a>0$, $\zeta(s,a)$ has a meromorphic continuation to the whole complex plane, with a single simple pole at $s=1$ of residue $1$. In the half-plane $\operatorname{Re}(s)>0$, $s\neq1$, one representation is
\[
 \zeta(s,a)=\frac{a^{1-s}}{s-1}+\frac12a^{-s}
 -s\int_0^{\infty}\frac{\widetilde B_1(x)}{(x+a)^{s+1}}\,dx,
 \qquad \widetilde B_1(x)=\{x\}-\frac12,
\]
where $\{x\}$ denotes the fractional part of $x$.
\end{theorem}

\begin{proof}
Euler's summation formula applied to $f(x)=(x+a)^{-s}$ gives the displayed representation first where the defining series converges. Since $\widetilde B_1$ is bounded and periodic with mean zero, its integral in the representation converges for $\operatorname{Re}(s)>0$ and defines a holomorphic function there. The only pole in that half-plane comes from $a^{1-s}/(s-1)$, which has residue $1$. Repeated Euler summation, using higher Bernoulli polynomials, continues the function to successively larger half-planes without introducing further poles. This gives the stated continuation.
\end{proof}

\begin{corollary}[Continuation of Dirichlet $L$-functions]
Every Dirichlet $L$-function has a meromorphic continuation to the whole complex plane. If $\chi$ is nonprincipal, $L(s,\chi)$ is entire. If $\chi$ is principal modulo $q$, then
\[
 L(s,\chi)=\zeta(s)\prod_{p\mid q}(1-p^{-s})
\]
and it has a simple pole at $s=1$.
\end{corollary}

\begin{proof}
Use the finite Hurwitz representation and its meromorphic continuation. Every Hurwitz zeta function has residue $1$ at $s=1$, so the possible residue in the finite sum is proportional to $\sum_{a=1}^{q}\chi(a)$. This sum is zero for nonprincipal $\chi$. For the principal character, removing the Euler factors corresponding to primes dividing $q$ gives the displayed formula.
\end{proof}

\begin{definition}[Gamma function]
For $\operatorname{Re}(s)>0$, the gamma function is
\[
 \Gamma(s)=\int_0^{\infty}e^{-t}t^{s-1}\,dt.
\]
It has a meromorphic continuation to the whole complex plane.
\end{definition}

\begin{theorem}[Functional equation for the Riemann zeta function]
The meromorphic continuation of $\zeta$ satisfies
\[
 \zeta(s)=2^s\pi^{s-1}\sin\!\left(\frac{\pi s}{2}\right)
 \Gamma(1-s)\zeta(1-s),
\]
where $\Gamma$ is the gamma function. The identity is understood meromorphically.
\end{theorem}

This functional equation relates values of $\zeta$ on opposite sides of the line $\operatorname{Re}(s)=1/2$ and is a central tool in the study of its zeros \cite[Chapter 12]{Apostol1976-xj}.

\subsection{Chebyshev Functions and Prime Distribution}

\begin{definition}[Chebyshev functions]
For $x\geq2$, define
\[
 \psi(x)=\sum_{n\leq x}\Lambda(n),
 \qquad
 \theta(x)=\sum_{p\leq x}\log p,
\]
where the second sum is over primes.
\end{definition}

\begin{proposition}[Prime-power decomposition]
For every $x\geq2$,
\[
 \psi(x)=\sum_{m=1}^{\lfloor\log_2 x\rfloor}\theta\!\left(x^{1/m}\right).
\]
\end{proposition}

\begin{proof}
In $\psi(x)$, only prime powers contribute. For each exponent $m\geq1$, the terms with $n=p^m\leq x$ contribute $\log p$, whose sum is $\theta(x^{1/m})$. No exponent greater than $\log_2x$ can occur.
\end{proof}

\begin{example}
The prime powers not exceeding $10$ are $2,4,8,3,9,5,$ and $7$. Therefore
\[
 \psi(10)=3\log2+2\log3+\log5+\log7=\log2520.
\]
\end{example}

\begin{theorem}[Chebyshev form of the prime number theorem]
The prime number theorem
\[
 \pi(x)\sim\frac{x}{\log x}\qquad(x\to\infty)
\]
 is equivalent to each of the asymptotic relations
\[
 \theta(x)\sim x
 \qquad\text{and}\qquad
 \psi(x)\sim x.
\]
Moreover, $\psi(x)-\theta(x)=O(\sqrt{x}\log^2x)$.
\end{theorem}

\begin{proof}
The relation between $\pi(x)$ and $\theta(x)$ follows by partial summation. The prime-power decomposition gives
\[
 0\leq\psi(x)-\theta(x)
 =\sum_{m=2}^{\lfloor\log_2x\rfloor}\theta(x^{1/m}).
\]
Using $\theta(y)\leq y\log y$ bounds this sum by $O(\sqrt{x}\log^2x)$, which is $o(x)$. Hence $\psi(x)\sim x$ if and only if $\theta(x)\sim x$, and both are equivalent to the prime number theorem.
\end{proof}

\subsection{A Contour Formula Involving the von Mangoldt Function}

Define the integrated Chebyshev function by
\[
 \psi_1(x)=\int_1^x\psi(u)\,du
 =\sum_{n\leq x}(x-n)\Lambda(n).
\]
The logarithmic derivative identity for $\zeta$ turns this weighted sum into a contour integral.

\begin{theorem}[Contour formula for $\psi_1$]
For $c>1$ and $x\geq1$,
\[
 \frac{\psi_1(x)}{x^2}
 =\frac{1}{2\pi i}\int_{c-i\infty}^{c+i\infty}
 \frac{x^{s-1}}{s(s+1)}\left(-\frac{\zeta'(s)}{\zeta(s)}\right)\,ds.
\]
The integral converges absolutely.
\end{theorem}

\begin{proof}
For $c>1$ and $y>0$, the residue theorem gives the Mellin integral
\[
 \frac{1}{2\pi i}\int_{c-i\infty}^{c+i\infty}
 \frac{y^s}{s(s+1)}\,ds
 =\begin{cases}
 1-y^{-1},&y\geq1,\\
 0,&0<y<1.
 \end{cases}
\]
Substitute $y=x/n$, multiply by $\Lambda(n)/x$, and sum over $n$. The identity $-\zeta'(s)/\zeta(s)=\sum_{n\geq1}\Lambda(n)n^{-s}$ and absolute convergence for $c>1$ permit interchanging the sum and integral. The resulting sum is
\[
 \frac{1}{x}\sum_{n\leq x}\Lambda(n)\left(1-\frac{n}{x}\right)
 =\frac{1}{x^2}\sum_{n\leq x}(x-n)\Lambda(n)
 =\frac{\psi_1(x)}{x^2},
\]
which proves the formula.
\end{proof}

This representation is one of the analytic links between the zeros and poles of $\zeta(s)$ and the distribution of prime powers. Apostol uses it as a starting point for the analytic proof of the prime number theorem \cite[Chapter 13]{Apostol1976-xj}.